\section{Transfer between proper and weak measurable colorings}
\label{sec:transfer}

We prove Theorem~\ref{thm:transfer} for every positive integer $k$.
Starting from an arbitrary proper coloring, we first average its
translations and rotations on the algebraic plane.  The resulting
probability space carries label indicators, but its translations need
not be strongly continuous in the usual plane topology.  We project those
indicators onto the functions whose translation orbits are continuous in $L^2$.
The projection must preserve nonnegativity and the sum of the label
coordinates; the next lemma establishes this by identifying it as
conditional expectation.  Haar rigidity then preserves the zero unit
correlations, and one measurable sample yields a weak coloring.
The converse uses density points and finite graph compactness.

\subsection{The continuous part of a stationary labeling system}

Retain the groups \(E,K,D\), the Borel subgroup \(C\subset D\), and
the parametrization \(j:\R^2\to C\) from
Definition~\ref{def:spectral-groups}.  In the following lemma, the topology
on \(E\) used to define continuity is its usual subspace topology in
\(\R^2\), although its measure-preserving action is initially only an
action of the countable discrete group.

\begin{lemma}[Continuous vectors and positivity]
\label{lem:palettes-continuous-factor}
Let a standard probability space carry a measure-preserving action of
the additive group \(E\), with Koopman operators \(T_z\).  The subspace
\[
 \mathcal H_c
 =\left\{v\in L^2:
       \lim_{\substack{z\to0\\z\in E}}\norm{T_zv-v}_2=0\right\}
\]
is \(L^2\) of a factor.  It is also the spectral subspace on \(C\)
for the representation of \(E_{\rm disc}\).
The operators on this subspace extend to a strongly continuous
unitary representation \(T_x\), \(x\in\R^2\), which preserves the positive
cone and the constant one.

If the system also has rotations \(V_u\), \(u\in K\), satisfying
\(V_uT_zV_u^{-1}=T_{uz}\), then \(\mathcal H_c\) and its orthogonal
projection are rotation invariant.
\end{lemma}

\begin{proof}
The subspace is closed.  Indeed, if \(v\) is close in \(L^2\) to a
continuous vector \(w\), then
\[
 \norm{T_zv-v}_2
 \le 2\norm{v-w}_2+\norm{T_zw-w}_2.
\]
It contains the constants and is invariant under the translations.

Apply the spectral theorem for the discrete abelian group \(E\), with
dual group \(D\).  For \(v\in L^2\), write \(\lambda_v\) for its finite
scalar spectral measure, so
\[
 \langle v,T_zv\rangle
   =\int_D\psi_z(d)\,d\lambda_v(d),\qquad
 \norm{T_zv-v}_2^2
   =\int_D\abs{\psi_z(d)-1}^2\,d\lambda_v(d).
\]
If \(\lambda_v(C^c)=0\), the second integral tends to
zero as \(z\to0\) in the usual topology, by dominated convergence and
the continuous-character parametrization.

Conversely, suppose \(v\in\mathcal H_c\).  Its coefficient
\(F_v(z)=\langle v,T_zv\rangle\) is uniformly continuous on \(E\):
\[
 \abs{F_v(z+h)-F_v(z)}
 \le \norm{v}_2\,\norm{T_hv-v}_2 .
\]
It extends to a continuous positive-definite function on \(\R^2\).
Positive definiteness follows by approximating each finite list of
real-plane points by points of \(E\) and passing to the limit in its
defining quadratic form.  Bochner's theorem supplies a finite measure
on \(\R^2\) representing this extension.  Its pushforward under \(j\)
has the same Fourier coefficients on \(D\) as \(\lambda_v\), so Fourier
uniqueness identifies the two measures.  Thus \(\lambda_v(C^c)=0\).
The representation result is the abelian spectral theorem
(also called the Stone--Naimark--Ambrose--Godement theorem); it and
Bochner's theorem are given in
\citet[Theorems~D.3.1 and~D.2.2]{Bekka2008}. The spectral theorem
is applied to the discrete locally compact group $E$, while
Bochner's theorem is applied to $\R^2$.  This identifies
\(\mathcal H_c\) with the range of the spectral projection on the Borel
set \(C\).

To obtain the factor assertion, it is essential that the original
\(T_z\) are Koopman operators.  They commute with complex conjugation
and with pointwise composition.  For bounded \(v,w\in\mathcal H_c\),
\[
 \norm{T_z(vw)-vw}_2
 \le \norm{v}_\infty\norm{T_zw-w}_2
     +\norm{w}_\infty\norm{T_zv-v}_2,
\]
so \(vw\in\mathcal H_c\).  For a Lipschitz function
\(\Phi:\C\to\C\),
\[
 \norm{T_z\Phi(v)-\Phi(v)}_2
 \le {\rm Lip}(\Phi)\norm{T_zv-v}_2 .
\]
In particular, radial projection of \(v\) onto a bounded complex disk
is still in \(\mathcal H_c\).  Such projections converge to \(v\)
in \(L^2\) as the disk radius tends to infinity.

Let \(\mathcal F_c\) be the completed factor generated by the bounded
members of \(\mathcal H_c\).  Those members form a bounded unital
star-algebra.  Its \(L^2\) closure is \(L^2(\mathcal F_c)\):
for a finite list of its bounded generators, polynomials in their
coordinates and conjugates approximate continuous functions on the
compact range by Stone--Weierstrass; continuous functions are dense
in \(L^2\) of the corresponding finite-dimensional distribution.
Finite-coordinate measurable functions are dense in the generated
factor, by the monotone-class theorem.  The closedness already proved
therefore gives \(L^2(\mathcal F_c)\subset\mathcal H_c\).
Conversely, the bounded truncations of every \(v\in\mathcal H_c\)
are \(\mathcal F_c\)-measurable and converge to \(v\) in \(L^2\).
Hence \(\mathcal H_c=L^2(\mathcal F_c)\).  Its orthogonal projection
is conditional expectation.

For \(x\in\R^2\), choose \(z_n\in E\) with \(z_n\to x\).
For \(v\in\mathcal H_c\),
\[
 \norm{T_{z_n}v-T_{z_m}v}_2
   =\norm{T_{z_n-z_m}v-v}_2\longrightarrow0 .
\]
Define \(T_xv\) by this limit.  The definition is independent of the
approximating sequence, and strong limits together with the original
group law give the group law on \(\R^2\), inverses, and strong
continuity.  The operators are unitary.  The positive cone in \(L^2\)
is closed, so \(T_x\) preserves positivity, and \(T_x1=1\).
No pointwise action of \(\R^2\) is needed for these conclusions.

Finally, for \(u\in K\), covariance gives
\[
 \norm{T_zV_uv-V_uv}_2
 =\norm{T_{u^{-1}z}v-v}_2\longrightarrow0 .
\]
Thus rotations preserve \(\mathcal H_c\); their inverses do as well,
so its orthogonal projection commutes with them.
\end{proof}

\begin{proposition}[From proper to weak measurable colorings]
\label{prop:palettes-forward}
Let \(k\) be a positive integer.  If the plane admits an unrestricted
proper \(k\)-coloring, then it admits a weak measurable \(k\)-coloring.
\end{proposition}

\begin{proof}
\emph{An invariant law on proper labelings.}
Restrict the given coloring to \(E\).  In the compact metrizable space
\(\{1,\ldots,k\}^E\), let \(\Omega\) be the closed set of proper
labelings of the unit-distance graph on \(E\).  It is nonempty.
The countable group \(E\rtimes K\) acts continuously on \(\Omega\);
translations and rotations preserve the unit-edge constraints.
This group is amenable, being an extension of two abelian groups
\citep[Theorem~G.2.1 and Proposition~G.2.2(ii)]{Bekka2008}.
The F{\o}lner criterion \citep[Theorem~G.5.1]{Bekka2008},
applied to successive finite subsets of this countable group, gives
a finite-set F{\o}lner sequence.
Averages of a point mass over a finite-set F{\o}lner sequence have a
weakly convergent subsequence, and every limit is invariant.  This
gives an invariant Borel probability \(\mu\) on \(\Omega\).

Let \(f_i\) be the indicator that the origin has label \(i\).
Choose the translation convention so \(T_zf_i\) is the indicator
of label \(i\) at \(z\).  The rotations about the origin fix each
\(f_i\), and properness gives
\begin{equation}\label{eq:palettes-proper-correlation}
 \int_\Omega f_i\,T_uf_i\,d\mu=0
 \qquad (u\in K,\ 1\le i\le k).
\end{equation}
\smallskip
\noindent\emph{The continuous factor and zero correlations.}
Use Lemma~\ref{lem:palettes-continuous-factor} to write
\[
 p_i=\mathbb E(f_i\mid\mathcal F_c),\qquad b_i=f_i-p_i.
\]
Then \(0\le p_i\le1\) and \(\sum_i p_i=1\).
Both \(p_i\) and \(b_i\) are rotation-fixed.  The scalar translation
spectral measure of \(b_i\) is \(K\)-invariant: covariance gives
equality of its Fourier coefficients at \(z\) and \(uz\), and
Fourier uniqueness gives invariance of the measure.  It gives \(C\) measure zero, by the spectral-subspace assertion of the lemma.
If it is nonzero, normalize it and apply Theorem~\ref{thm:rigidity}.
The scalar spectral measure is therefore a multiple of Haar measure,
including the zero multiple.  It follows that
\[
 \langle b_i,T_zb_i\rangle=0\qquad(z\in E\setminus\{0\}).
\]
The subspace \(\mathcal H_c\) and its orthogonal complement are
translation invariant, so both cross terms in
\eqref{eq:palettes-proper-correlation} vanish.  Consequently
\[
 \int_\Omega p_i\,T_up_i\,d\mu=0\qquad(u\in K).
\]
The algebraic unit directions are dense in the usual unit circle,
as is seen already from the rational tangent-half-angle parametrization.
Strong continuity of the extended operators therefore gives this
identity for every usual unit direction \(u\in S^1\).
For every \(x\in\R^2\), unitarity and the group law then give
\begin{equation}\label{eq:palettes-field-correlation}
 \int_\Omega (T_xp_i)(T_{x+u}p_i)\,d\mu
 =\langle T_xp_i,T_xT_up_i\rangle
 =\langle p_i,T_up_i\rangle=0 .
\end{equation}
The functions here are nonnegative equivalence classes.
This calculation uses no pointwise realization or multiplicativity
of the extended \(\R^2\) action.

\smallskip
\noindent\emph{A measurable sample.}
The continuous map
\(x\mapsto T_xp_i\) takes values in the separable space \(L^2(\mu)\).
On a bounded spatial domain, approximate it by finite-valued simple
maps in the \(L^2(\mu)\) norm, choosing the integrals of their squared
approximation errors to be summable.  Choose measurable representatives
for their finitely many values.  The resulting jointly measurable
functions converge in the product \(L^2\) space.  A subsequence converges
in \(L^2(\mu)\) for almost every \(x\), so its product-space limit represents
\(T_xp_i\) at almost every \(x\).  Applying this on a countable
exhaustion, or on its disjoint measurable annuli, gives a jointly
measurable field \(H_i(x,\omega)\) on \(\R^2\times\Omega\), correct
as an \(L^2(\mu)\) representative for almost every \(x\).
Since the probability model is standard, Borel versions may be used.
Taking real parts and clamping each \(H_i\) to \([0,1]\) preserves this representation
property and makes all the fields everywhere nonnegative.

There is a plane-null set \(N\) outside which all these representations
are correct.  For almost every \(x\), their sum is one for
\(\mu\)-almost every \(\omega\).  Moreover, on every bounded
spatial domain,
\[
 \int\!\int_{S^1}
   \bigl(\ind_N(x)+\ind_N(x+u)\bigr)\,d\sigma(u)\,dx=0.
\]
The assertion about \(x+u\) follows by translation invariance of
Lebesgue measure for each fixed \(u\), followed by Fubini.
Thus \eqref{eq:palettes-field-correlation} applies to the chosen
representatives for almost every pair \((x,u)\).
Tonelli gives, for every positive integer \(m\),
\[
 \int_\Omega\sum_{i=1}^k
       \int_{B(0,m)}\int_{S^1}
       H_i(x,\omega)H_i(x+u,\omega)\,d\sigma(u)\,dx\,d\mu(\omega)=0.
\]
Choose one sample \(\omega\) for which all the inner integrals vanish,
and for which \(\sum_iH_i(x,\omega)=1\) for almost every \(x\).
There are only countably many conditions here.  At every point
where the sum is one, assign the least index with positive field;
on the remaining plane-null set assign any index.  This is measurable.
A same-color pair outside that null set would have positive product
of the corresponding fields, which occurs only on a
\(dx\,d\sigma(u)\)-null set.  Hitting the spatial exceptional set at
either endpoint is also a product-null event.  This proves the
weak measurable condition.
\end{proof}

\subsection{Density points and the converse transfer}

The next argument is a version of the density-point method used by
\citet{Falconer1981} and \citet[Lemma~1 and Proposition~1]{Payne2009}.
Here the hypothesis excludes unit pairs only in product measure, so we
first show directly that density-one points still form independent sets.

\begin{lemma}[Density-point properness]\label{lem:palettes-density-proper}
In a weak measurable \(k\)-coloring with classes \(A_1,\ldots,A_k\),
the full-density-point set of each \(A_i\) contains no unit pair.
These density-point sets are disjoint and their union is conull.
\end{lemma}

\begin{proof}
Suppose \(x\) and \(y=x+u\), with \(\abs{u}=1\), both have density one
of \(A_i\).  For \(h>0\) small, choose \(x'\) uniformly in \(B(x,h)\)
and independently choose \(u'\) uniformly on an arc of positive
angular measure satisfying \(\abs{u'-u}<h\).
Density at \(x\) gives
\(\mathbb P(x'\notin A_i)=o(1)\) as \(h\downarrow0\).
The distribution of \(x'+u'\) is a mixture of translated uniform
disk distributions.  Its density is bounded by \(1/(\pi h^2)\),
and its support is in \(B(y,2h)\).  Therefore
\[
 \mathbb P(x'+u'\notin A_i)
 \le \frac{\abs{B(y,2h)\setminus A_i}}{\pi h^2}=o(1).
\]
For small \(h\) the probability that both points lie in \(A_i\)
is positive.  But the pair-parameter distribution is absolutely
continuous with respect to \(dx\,d\sigma(u)\), and the two points
have distance one.  This contradicts the weak condition.

Lebesgue differentiation makes the union of the density-point sets
conull.  Two different color classes, being disjoint up to null sets,
cannot both have density one at the same point, so these sets are
disjoint.
\end{proof}

\begin{proof}[Proof of Theorem~\ref{thm:transfer}]
The forward implication is Proposition~\ref{prop:palettes-forward}.
For the reverse implication, Lemma~\ref{lem:palettes-density-proper}
gives a conull set partitioned into strictly unit-independent
density-point classes.  Any fixed finite configuration of plane
points has a translate entirely inside that conull set: the permissible
translation vectors are a finite intersection of conull sets.
Its unit-distance graph therefore has a proper \(k\)-coloring.

In the compact product space \(\{1,\ldots,k\}^{\R^2}\), each unit-edge
condition is closed and depends on two coordinates.  Every finite
family of these conditions involves a finite configuration and
is satisfiable by the preceding paragraph.  The finite intersection
property, using product compactness in ZFC, produces a coloring
satisfying every unit-edge condition.  This is the graph compactness
principle of \cite[Theorem 1]{deBruijnErdos1951}; the argument just
given specifies its application to the unrestricted plane.
It proves \eqref{eq:source-1}.
\end{proof}

